
\documentclass{article}
\usepackage{colortbl}
\usepackage{makecell}
\usepackage{multirow}
\usepackage{supertabular}

\begin{document}

\newcounter{utterance}

\centering \large Interaction Transcript for game `cladder', experiment `full\_v1.5\_default', episode 8179 with BSU{-}SLIM\_\_Qwen3.5{-}9B{-}Base\_\_playpen{-}prm{-}9b{-}best\_of\_n.
\vspace{24pt}

{ \footnotesize  \setcounter{utterance}{1}
\setlength{\tabcolsep}{0pt}
\begin{supertabular}{c@{$\;$}|p{.15\linewidth}@{}p{.15\linewidth}p{.15\linewidth}p{.15\linewidth}p{.15\linewidth}p{.15\linewidth}}
    \# & \multicolumn{2}{c}{Player} && \multicolumn{2}{c}{Game Master} \\
    \hline

    \theutterance \stepcounter{utterance}  
    & & & \multicolumn{4}{p{0.6\linewidth}}{
        \cellcolor[rgb]{0.9,0.9,0.9}{
            \makecell[{{p{\linewidth}}}]{
                \texttt{\tiny{[P1$\langle$GM]}}
                \texttt{You are an expert in causal inference. The following question is not a typical commonsense query, but rather a meticulously designed question created by a professor specializing in causal inference, intended to assess the students' mastery of the course content.} \\
\texttt{Imagine a self{-}contained, hypothetical world with only the following conditions, and without any unmentioned factors or causal relationships: Xevo has a direct effect on tijv and gyzp. Tijw has a direct effect on tijv and gyzp. Tijv has a direct effect on gyzp. Tijw is unobserved. Method 1: We look directly at how xevo correlates with gyzp in general. Method 2: We look at this correlation case by case according to tijv. To understand how xevo affects gyzp, is it more correct to use the Method 1 than Method 2?} \\
\texttt{Start your answer with 'Yes' or 'No', followed by additional reasoning or evidence to support your explanation.} \\
            }
        }
    }
     \\ \\

    \theutterance \stepcounter{utterance}  
    & \multicolumn{4}{p{0.6\linewidth}}{
        \cellcolor[rgb]{0.561,0.737,0.561}{
            \makecell[{{p{\linewidth}}}]{
                \texttt{\tiny{[P1$\rangle$GM]}}
                \texttt{No} \\
\\ 
\texttt{In this causal world, Method 1 is incorrect because it suffers from **omitted variable bias** due to the unobserved confounder \$tijw\$. Since \$tijw\$ has a direct effect on both the exposure (\$xevo\$) and the outcome (\$gyzp\$), \$xevo\$ and \$gyzp\$ are correlated even if \$xevo\$ has no direct causal effect on \$gyzp\$, purely because they are both driven by \$tijw\$. Method 1, which simply looks at the marginal correlation between \$xevo\$ and \$gyzp\$, conflates the direct effect of \$xevo\$ with the spurious correlation induced by \$tijw\$.} \\
\\ 
\texttt{Method 2 is the correct approach because it conditions on \$tijv\$. Since \$xevo\$ has a direct effect on \$tijv\$, and \$tijv\$ has a direct effect on \$gyzp\$, \$tijv\$ acts as a **mediator** for the \$xevo \textbackslash{}rightarrow gyzp\$ pathway. By conditioning on the mediator (\$tijv\$), the correlation between \$xevo\$ and \$gyzp\$ in this stratified analysis isolates the direct path from \$xevo\$ to \$gyzp\$ that does not go through \$tijv\$. Crucially, because \$tijw\$ only affects \$gyzp\$ (and not \$tijv\$ based on the provided graph where \$xevo\$ affects \$tijv\$) and \$tijw\$ is unobserved, conditioning on \$tijv\$ does *not* open a backdoor path through \$tijw\$. Therefore, the conditional correlation in Method 2 provides an unbiased estimate of the direct causal effect of \$xevo\$ on \$gyzp\$, whereas Method 1 provides a biased estimate of the total effect obscured by the confounding influence of \$tijw\$.} \\
            }
        }
    }
    & & \\ \\

    \theutterance \stepcounter{utterance}  
    & & & \multicolumn{2}{p{0.3\linewidth}}{
        \cellcolor[rgb]{0.9,0.9,0.9}{
            \makecell[{{p{\linewidth}}}]{
                \texttt{\tiny{[GM$|$GM]}}
                \texttt{No} \\
            }
        }
    }
    & & \\ \\

    \theutterance \stepcounter{utterance}  
    & & & \multicolumn{2}{p{0.3\linewidth}}{
        \cellcolor[rgb]{0.9,0.9,0.9}{
            \makecell[{{p{\linewidth}}}]{
                \texttt{\tiny{[GM$|$GM]}}
                \texttt{yes} \\
            }
        }
    }
    & & \\ \\

    \theutterance \stepcounter{utterance}  
    & & & \multicolumn{2}{p{0.3\linewidth}}{
        \cellcolor[rgb]{0.9,0.9,0.9}{
            \makecell[{{p{\linewidth}}}]{
                \texttt{\tiny{[GM$|$GM]}}
                \texttt{game\_result = LOSE} \\
            }
        }
    }
    & & \\ \\

\end{supertabular}
}

\end{document}
